\section{Solvent and product use}
\label{sec:solvents}

\subsection{Scope and inventory mapping}

This block targets the area-source component of GNFR~E, which groups
solvent and product use activities under the broad heading of product
application, surface treatment, degreasing, domestic solvent use, and
related diffuse processes. In contrast to public power, these emissions
are not attached to a small number of explicit stacks, and in contrast
to residential combustion they are not driven primarily by one dominant
activity variable such as heat demand. The spatial problem is therefore
to redistribute coarse CAMS area-source mass toward fine-grid locations
that are plausible for the different types of solvent use aggregated
inside GNFR~E.

The implementation is restricted to the CAMS area-source geometry
within GNFR~E. For each CAMS parent cell and pollutant, the final
fine-grid weights are normalised so that the redistributed pixel shares
sum to unity inside that parent cell. The output is thus a
mass-conserving multi-band weight raster, one band per pollutant, ready
to be combined with the corresponding CAMS cell totals.

\subsection{Conceptual design}

GNFR~E is internally heterogeneous. Domestic product use, road paving,
roofing, coating, printing, chemical products, dry cleaning, and
industrial degreasing do not share the same fine-scale geography.
A single land-cover mask would therefore reproduce the same weakness
identified earlier for agriculture and residential combustion: one
sector-level proxy would force all solvent-related processes to follow a
single spatial pattern.

The upgraded design separates the problem into three layers:
\begin{itemize}
  \item a set of fine-grid indicator layers describing settlement,
        services, industry, and transport-related infrastructure;
  \item an intermediate set of four archetype activity surfaces
        representing household, service, industrial, and infrastructure
        solvent-use environments;
  \item a final set of nine solvent subsectors, each expressed as a
        mixture of the four archetypes and then recombined with
        pollutant-specific national shares.
\end{itemize}

This creates a two-stage weighting system. First, spatial indicators
describe where each type of solvent-related activity is likely to occur.
Second, national subsector shares determine how much of each pollutant
should follow each solvent-use subsector. The resulting method remains
mass conserving within each CAMS parent cell while allowing pollutant
bands to differ spatially.

The archetype definitions and subsector mixing coefficients are
summarised in Appendix~\ref{app:solvents_downscaling},
Tables~\ref{tab:solvents_archetypes} and
\ref{tab:solvents_subsectors_beta}.

\subsection{Fine-grid activity surfaces}

Let $g$ denote a fine pixel on the national reference grid. The method
first constructs a set of non-negative indicator layers $z_{g,u}$ from
harmonised ancillary datasets. These indicators represent, at fine
resolution, the main spatial environments in which solvent use is
expected to occur:
\begin{itemize}
  \item population and residential urban fabric for household-oriented
        uses;
  \item service-oriented urban and commercial land for retail and
        tertiary activities;
  \item industrial land and building footprints for manufacturing and
        product-processing uses;
  \item road length, road hierarchy, transport-related land, and built
        surfaces for infrastructure-linked activities.
\end{itemize}

The underlying indicators combine gridded population, CORINE land-cover
classes, and OpenStreetMap-derived roads and land-use polygons. Each
indicator is first transformed to a unitless relevance field
$\tilde{z}_{g,u} \in [0,1]$ using a clipped quantile scaling:
\begin{equation}
  \tilde{z}_{g,u} =
  \mathrm{clip}_{[0,1]}
  \left(
    \frac{
      \min\!\left(\max(z_{g,u}, q_u^{0.01}), q_u^{0.99}\right) - q_u^{0.01}
    }{
      q_u^{0.99} - q_u^{0.01}
    }
  \right),
\end{equation}
where $q_u^{0.01}$ and $q_u^{0.99}$ are the 1st and 99th percentiles of
positive finite values for indicator $u$. This scaling preserves
relative spatial contrasts while preventing a small number of extreme
pixels from dominating the proxy.

\begin{table}[htbp]
  \centering
  \caption{Spatial evidence used for the solvent archetype proxy surfaces.}
  \label{tab:solvents-spatial-rules}
  \scriptsize
  \begin{tabular}{p{0.18\textwidth} p{0.28\textwidth} p{0.42\textwidth}}
    \toprule
    Archetype & Spatial weighting & Retained evidence \\
    \midrule
    Household &
    $\lambda_{\mathrm{pop}}=0.70$, $\lambda_{\mathrm{res}}=0.30$. &
    Gridded population and residential urban fabric classes. \\
    Service &
    $\lambda_{\mathrm{urban}}=0.40$, $\lambda_{\mathrm{service\,land}}=0.35$, $\lambda_{\mathrm{OSM}}=0.25$. &
    Urban fabric, commercial and service land, and OSM retail, commercial, civic, and public land-use features. \\
    Industrial &
    $\lambda_{\mathrm{CLC}}=0.35$, $\lambda_{\mathrm{OSM}}=0.35$, $\lambda_{\mathrm{buildings}}=0.30$. &
    Industrial and construction-related land, OSM industrial, commercial, warehouse, service, and aeroway features, and building footprints. \\
    Infrastructure &
    $\lambda_{\mathrm{road}}=0.35$, $\lambda_{\mathrm{weighted\,road}}=0.35$, $\lambda_{\mathrm{transport}}=0.20$, $\lambda_{\mathrm{roof}}=0.10$. &
    Road length, hierarchy-weighted road length, transport-related land, and roof-area support. \\
    \bottomrule
  \end{tabular}
\end{table}

The normalised indicators are then aggregated into four archetype
surfaces, indexed by $a \in \{\mathrm{house}, \mathrm{serv},
\mathrm{ind}, \mathrm{infra}\}$:
\begin{equation}
  A_g^{(a)} =
  \sum_{u \in \mathcal{U}(a)}
  \lambda_u^{(a)} \, \tilde{z}_{g,u},
\end{equation}
with weights $\lambda_u^{(a)}$ renormalised over the indicators actually
available for that archetype. The four archetypes may be interpreted as
first-order activity surfaces for household solvent use, service-sector
uses, industrial uses, and infrastructure-linked uses respectively.

\subsection{Subsector activity and pollutant shares}

The solvent block is decomposed into nine subsectors corresponding to
the CEIP 2.D.3 subdivision:
\begin{equation}
  s \in \{
  \mathrm{2D3a},\mathrm{2D3b},\ldots,\mathrm{2D3i}
  \}.
\end{equation}
These subsectors cover domestic solvent use, road paving, asphalt
roofing, coating, degreasing, dry cleaning, chemical products,
printing, and a residual ``other'' category.

Each subsector is represented as a convex mixture of the four archetype
surfaces:
\begin{equation}
  R_g^{(s)} =
  \sum_{a}
  \beta_{s,a} \, A_g^{(a)},
\end{equation}
where the coefficients $\beta_{s,a}$ satisfy
\begin{equation}
  \sum_a \beta_{s,a} = 1
  \qquad \forall s.
\end{equation}
The $\beta$ matrix expresses the substantive interpretation of each
solvent-use subsector. For example, road paving is mapped entirely to
the infrastructure archetype, dry cleaning to the service archetype,
chemical products to the industrial archetype, and mixed activities such
as coating or ``other'' to weighted combinations of several archetypes.

Pollutant-specific national shares are derived from national 2.D.3
totals. Let $E_{r,s,p}^{\mathrm{nat}}$ denote the reported national
emission total for country $r$, subsector $s$, and pollutant $p$. The
subsector share is
\begin{equation}
  \alpha_{r,s,p} =
  \frac{E_{r,s,p}^{\mathrm{nat}}}
       {\displaystyle\sum_{s'} E_{r,s',p}^{\mathrm{nat}}}.
  \label{eq:solvents-alpha}
\end{equation}
Two completion rules are applied to make these ratios operational for
all pollutants:
\begin{itemize}
  \item if a pollutant is reported only for a subset of subsectors, the
        non-reported subsectors retain zero and the reported subsectors
        are normalised to sum to one;
  \item if a pollutant has no solvent subsector information at all, the
        nine subsectors receive equal shares, that is
        $\alpha_{r,s,p}=1/9$ for all $s$.
\end{itemize}

This choice reflects the intended interpretation of the CEIP data. A
missing pollutant row does not imply the absence of GNFR~E emissions in
the CAMS inventory; it only means that no national solvent subsector
split is available for that pollutant. In that case, an uninformative
equal split is preferred to borrowing the mixture of another pollutant.

\subsection{Within-cell normalisation and final weights}

Let $m$ denote a CAMS GNFR~E area-source parent cell and let
$g \in m$ denote fine pixels whose centres fall inside that parent cell.
Each raw subsector surface $R_g^{(s)}$ is normalised within the
corresponding CAMS parent cell:
\begin{equation}
  \rho_{g \mid m}^{(s)} =
  \frac{R_g^{(s)}}
       {\displaystyle\sum_{g' \in m} R_{g'}^{(s)}}.
\end{equation}
If the denominator is zero for a given subsector and CAMS cell, a
generic solvent proxy is used instead:
\begin{equation}
  G_g = \sum_a \eta_a A_g^{(a)},
\end{equation}
where $\eta_a$ are equal archetype weights. If that generic proxy is
also zero within the CAMS cell, the algorithm falls back to a uniform
distribution over the valid pixel support of the cell. This ensures that
no CAMS area-source mass is lost simply because a particular fine-grid
indicator is absent locally.

The final pollutant-specific weight in pixel $g$ for CAMS cell $m$ is
then
\begin{equation}
  W_{g,m,p}^{\mathrm{solv}} =
  \sum_s \alpha_{r,s,p} \, \rho_{g \mid m}^{(s)}.
\end{equation}
By construction,
\begin{equation}
  \sum_{g \in m} W_{g,m,p}^{\mathrm{solv}} = 1
  \qquad \forall m,p,
\end{equation}
so the fine-grid redistribution is mass conserving inside every CAMS
parent cell.

If $E_{m,p}^{\mathrm{CAMS}}$ denotes the CAMS GNFR~E area emission for
parent cell $m$ and pollutant $p$, the gridded allocation is
\begin{equation}
  e_{g,m,p}^{\mathrm{solv}} =
  E_{m,p}^{\mathrm{CAMS}} \,
  W_{g,m,p}^{\mathrm{solv}}.
\end{equation}
The method therefore separates three distinct tasks: activity
representation through archetypes, pollutant-specific subsector shares
through $\alpha$, and mass-conserving redistribution within CAMS cells
through $\rho$.

\subsection{Interpretation of activity and alpha ratios}

The activity term in this block should be understood as a spatial
likelihood surface rather than as a direct estimate of solvent use in
physical units. Population density, service land, industrial land,
building surfaces, and road infrastructure are not themselves emissions,
but they are informative about where the different solvent-use
subsectors are likely to occur. The four archetypes therefore play the
same conceptual role as source-specific proxy families in the other
chapter sections: they encode relative spatial relevance.

The $\alpha$ ratios have a different role. They do not alter the
fine-scale geography directly; instead, they determine how much each
pollutant draws from each solvent subsector. A pollutant associated
mainly with domestic product use will inherit more of the household
archetype, while a pollutant concentrated in chemical products or
degreasing will inherit more of the industrial mixture. Spatial
differences between pollutant bands thus arise from the interaction
between the common fine-grid activity surfaces and the
pollutant-specific national subsector shares.

\subsection{Limitations}

The method remains a first-order allocation scheme for a highly diverse
sector. Several limitations follow directly from that role.

First, the archetype surfaces are inferred from generic spatial
indicators rather than observed solvent-use statistics. They distinguish
plausible environments of use, but they do not recover product-specific
consumption or facility-level activity directly.

Second, the national subsector shares are only as detailed as the
available 2.D.3 reporting. Where a pollutant is reported for only a few
subsectors, the method preserves that selectivity. Where no solvent
subsector split is reported, the equal-share fallback is intentionally
uninformative and should be read as a regularisation device rather than
as evidence of equal real-world contributions.

Third, the infrastructure-related component relies partly on
OpenStreetMap geometry. Although this is a major improvement over a pure
land-cover proxy, completeness and tagging quality may still vary across
space.

Despite these limitations, the GNFR~E workflow represents a substantial
advance over a single composite proxy. It allows solvent emissions to
respond differently to household, service, industrial, and
infrastructure environments while preserving pollutant-specific national
subsector structure and exact mass conservation within each CAMS parent
cell.
